\documentclass{book}
\usepackage{hyperref}
\usepackage[margin=1in]{geometry}
\usepackage{ulem}
\usepackage{tipa} % Add this line to support \textipa command
\hypersetup{
    colorlinks=true,
    linkcolor=black,
    filecolor=magenta,      
    urlcolor=cyan,
    }

\urlstyle{same}

\title{French Conjugation}
\author{Junzhang Luo}

\begin{document}

\tableofcontents

\section{Day 1}
    Les jours de la semaine, je retourne chez moi autout dix-six ou dix-sept heures.
    Tout d'abord, avant de dîner, je me douche et s'habille le tenue de la maison.
    En géneral, je mange un bol de riz avec un ou deux plats. 
    Après le dîner, je consacre des temps à l'apprentissage de français en dévelopant les logciels.
    J'essaie toujours de apprendre nouveaux quelque chose.
    \\
    Pour le week-end, je apprends de français, et parfois, je sortis à reconnetre avec mes amis et on fait le sports, tel que jouer au badminton, au basket-ball et jeux de vidéo.

    \subsection{Correction}
        Les jours de la semaine, je rentre chez moi vers seize ou dix-sept heures.
        Tout d'abord, avant de dîner, je me douche et je m'habille avec des vêtements confortables pour la maison.
        En général, je mange un bol de riz avec un ou deux plats. 
        Après le dîner, je consacre du temps à l'apprentissage du français tout en développant des logiciels.
        J'essaie toujours d'apprendre quelque chose de nouveau pour améliorer mes compétences.
        Le soir, je lis parfois un peu ou j'écoute du français pour progresser davantage.

    \subsection{C1}
        En semaine, je rentre généralement chez moi vers seize ou dix-sept heures. 
        Avant de dîner, je prends d'abord une douche et j'enfile des vêtements plus confortables.
        En général, mon repas se compose d'un bol de riz accompagné d'un ou deux plats. 
        Après le dîner, je consacre du temps à l'étude du français tout en travaillant sur le développement de logiciels.
        J'essaie toujours d'apprendre quelque chose de nouveau afin de progresser.
        Quand j'ai encore un peu d'énergie, j'écoute aussi du contenu en français pour m'habituer davantage à la langue.

    \begin{itemize}
        \item autour de / vers, use vers for time
        \item use definite article for a school subject or language
        \item faire du sport (fixed expression)
        \item jeux vidéo (fixed noun phrase)
    \end{itemize}

\section{Day 2}
    Je m'expatrie à Toronto dans 2017 et je habite ici depuis.
    Il y a pleusieurs des lieux d'intérêt dans la région de Toronto.
    En général, les gens jeunes traînent dans la centre ville Toronto, 
    parce qu'il y a le café, la barre, la centre commercial et les activités qui convenir à les gens jeunes.
    En revanche, les autres zones dehors la centre ville sont très agréable et beaux.
    Il y a beaucoup des parcs dispersée autour de la zone de Toronto, appropriée pour se promenade.
    On trouve l'île de centre Lac Ontario et c'est un lieux pacifique pour rester et apprécier de la nature.

    \subsection{Version Corrigée}
        Je me suis installé à Toronto en 2017 et j'y habite depuis. 
        Il y a plusieurs lieux d'intérêt dans la région de Toronto.
        En général, les jeunes aiment passer du temps dans le centre-ville de Toronto, parce qu'on y trouve des cafés, des bars, des centres commerciaux et diverses activités qui leur conviennent.
        En revanche, les quartiers situés en dehors du centre-ville sont très agréables et beaux.
        Il y a aussi beaucoup de parcs dispersés dans la région, qui sont parfaits pour se promener.
        On trouve également les îles de Toronto, situées sur le lac Ontario, et c'est un endroit paisble pour se détendre et profiter de la nature.

    \subsection{C1}
        Je me suis installé à Toronto en 2017 et j'y vis depuis. 
        Cette ville offre de nombreux lieux d'intérêt.
        En général, les jeunes aiment fréquenter le centre-ville,
        où l'on trouve une grande variété de cafés, de bars, de centres commerciaux et d'activités de loisirs.
        En revanche, les quartiers plus éloignés du centre se distinguent par leur calme et leur cadre agréab;e.
        Toronto compte également de nombreux parcs répartis dans toute la région, ce qui en fait un endroit idéal pour se promener.
        On peut aussi vister les îles de Toronto, sur le lac Ontario, qui constituent un lieu paisible où l'on peut se reposer et profiter de la nature.

    \subsection{Remarks}
        \begin{itemize}
            \item for a year, use ``en 2017'', instead of dans
            \item j'habite \dots depuis \dots, je me suis installé \dots, better than s'expatrier
            \item plusieurs, no de
            \item les jeunes, no need for ``gens''
            \item le centre-ville de
            \item when listing general places, use plural and usually des ``des''
            \item convenir à (to suit)
            \item en dehors de (outside of)
            \item beaucoup de for general quantity (not ``des'')
            \item after pour, use infinitive
            \item ``paisible'' is more natural for a place
        \end{itemize}

\section{Day 3}
    Je suis Junzhang et j'ai dix-huit ans. J'aime la musique, le sports et le science.
    J'ai commencé apprendre à jouer du piano quand j'avais cinq ans.
    J'ai également appris l'écriture caligraphique depuis six ans.
    Je suis allé l'école primaire à Beijing en Chine et j'ai m'installé au Canada en 2017.
    J'ai fais partie de la orchestre d'harmonie au lycée comme un joueur de la saxophone.
    Entre-temps, j'ai fais partie de l'ensemble de la saxophone.
    J'aime voyager aux lieux différentes et expérience les culture, goûter les repas locaux, interagir avec les gens.

    \subsection{Correction}
        Je suis Junzhang et j'ai dix-huit ans. 
        J'aime la musique, le sport et la science.
        J'ai commencé à jouer du piano quand j'avais cinq ans.
        J'apprends également la calligraphie depuis six ans.
        Je suis allé à l'école primaire à Beijing, en Chine, et je me suis installé au Canada en 2017.
        J'ai fait partie de l'orchestre d'harmonie au lycée en tant que saxophoniste.
        Entre-temps, j'ai aussi fait partie d'un ensemble de saxophones.
        J'aime voyager dans différents endroits, découvir différent cultures, goûter des plats locaux et interagir avec les gens.

    \subsection{Remarks}
        \begin{itemize}
            \item commencé à
            \item je suis allé à
            \item s'installer is reflexive, so takes être
            \item en tant que (more natural)
            \item dans différents endroits
            \item échanger (to exchange)
            \item échanger avec (talk with, have a conversation, share ideas)
            \item les habitants (local people)
        \end{itemize}

    \subsection{Adv}
        Je m'appelle Junzhang, j'ai dix-huit ans et je suis passionné par le musique, le sport et la science.
        Je joue du piano depuis l'âge de cunq ans et je pratique aussi la calligraphie depuis six ans.
        Né à Beijing, où j'ai fait mes études primaires, je me suis ensuite installé au Canada en 2018.
        Au lycée, j'ai fait partie de l'orchestre d'harmonie ainsi que d'un ensemble de saxophones.
        J'aime également voyager, découvrir d'autres cultures, goûter des spécialitiés locales et écahnger avec les habitants,
        car ces expériences me permettent d'élargir ma vison du monde.


\end{document}